% !TeX program = xelatex
% !TEX encoding = UTF-8
%
% Шаблон статьи «Журнала Прикладной Лженауки».
% Порядок разделов и обязательные элементы — см. guide/…guide.md, §2.
% Скопируйте этот файл под именем <slug>.tex, заполните TODO, соберите
% `cd articles && latexmk`, затем добавьте запись в articles/articles.json
% и прогоните `python scripts/gen_site.py`.

\documentclass[11pt,a4paper]{article}
\usepackage{zhpl}

% TODO: заголовок в два уровня — интригующий + строгий академический через двоеточие
\title{\textbf{TODO — интригующий заголовок:\\ TODO — строгий академический подзаголовок}}
\author{\zhplauthors}

% TODO: номер выпуска и DOI вида 10.9999/жпл.20XX.000N
\zhplsetup{TODO}{10.9999/жпл.2026.000TODO}

\begin{document}
\maketitle

\zhplbadges

% TODO: DOI и диапазон страниц — как в колонтитуле
\zhplciteas{Александров Р.\,Г., Антропикович К.\,С. (2026). TODO — краткое название.
  \textit{Журнал Прикладной Лженауки}, \textbf{1}(TODO), 1--TODO. DOI: 10.9999/жпл.2026.000TODO}

% TODO: графический реферат — 3–4 бокса со стрелками, логика статьи одним взглядом
\zhplgraphicabstract{%
  \begin{tikzpicture}[node distance=0.9cm, every node/.style={font=\footnotesize}]
    \node[draw,rounded corners,fill=blue!5,minimum width=2.4cm,minimum height=1cm,align=center] (n1) {TODO};
    \node[draw,rounded corners,fill=blue!5,minimum width=2.4cm,minimum height=1cm,align=center,right=of n1] (n2) {TODO};
    \node[draw,rounded corners,fill=green!8,minimum width=2.4cm,minimum height=1cm,align=center,right=of n2] (n3) {TODO};
    \draw[-{Latex[length=2.5mm]},very thick] (n1) -- (n2);
    \draw[-{Latex[length=2.5mm]},very thick] (n2) -- (n3);
  \end{tikzpicture}}{%
  TODO — одна фраза: от бытовой постановки к формальному инварианту, ценой всей информации.}

\begin{zhplhighlights}
\begin{itemize}
\setlength\itemsep{0pt}
\setlength\parskip{0pt}
\item TODO — 3–5 пунктов, каждый со ссылкой на Теорему N.
\end{itemize}
\end{zhplhighlights}

\begin{zhplreviews}
\textit{«TODO — переубеждённый скептик.»} -- Рецензент 1\\
\textit{«Замечаний нет.»} -- Рецензент 2\\
\textit{«TODO — мета-шутка: согласно Теореме N этой же статьи объективная рецензия невозможна.»} -- Рецензент 3
\end{zhplreviews}

\begin{zhplsignificance}
TODO — один грандиозный абзац в духе PNAS: наивная (не-строгая) версия темы веками
считалась неразрешимой; настоящая работа впервые показывает обратное; результат
применим ко всей человеческой цивилизации.
\end{zhplsignificance}

\begin{abstract}
\noindent Здесь мы впервые показываем, что TODO --- не <<бытовая>>, а строгая
TODO-закономерность. Далее: проблема $\to$ метод $\to$ результаты $\to$ ограничения.

\medskip
\noindent\textbf{Ключевые слова:} TODO, TODO, лженаука
\end{abstract}

\section{Введение}
% TODO: постановка проблемы как «на самом деле геометрической/термодинамической и т.п.,
% а не той, за которую её принимали».

\section{TODO — раздел 1: формальный объект}
% TODO: ввести метрику/тензор/оператор/потенциал по теме статьи.

\begin{definition}[TODO — говорящее название]
TODO — строгая формулировка.
\end{definition}

\begin{figure}[htbp]
\centering
% TODO: иллюстрация, посчитанная точно (pgfplots с настоящей формулой; касательные —
% по вычисленной производной; сфера — полноценной сеткой). Ничего «на глаз». guide §3.
\caption{TODO. Ссылка в тексте через \ref{fig:todo1}.}
\label{fig:todo1}
\end{figure}

\section{TODO — раздел 2}
\begin{theorem}[TODO — говорящее название]
TODO — строгая формулировка.
\end{theorem}

\section{TODO — раздел 3}
% TODO: ещё 1–2 содержательных раздела, в каждом теорема/определение + иллюстрация.

\section{Обсуждение и ограничения}
% TODO: 2–3 пункта, честно признающие, почему метод не работает (потеря информации
% при свёртке, локальность, неустранимая корреляция). Финал: «это можно рассматривать
% как единственный строгий результат работы».

\section{Заключение}
% TODO: коротко; последняя фраза — про «достаточное основание для дальнейшего
% финансирования» или аналог.

\section*{Вклад авторов}
Р.\,Г.~Александров: концептуализация, методология, написание -- черновик.
К.\,С.~Антропикович: программное обеспечение, визуализация, формальный анализ,
написание -- рецензирование и редактирование, техническая поддержка бесконечного терпения.

\section*{Финансирование}
Исследование выполнено при поддержке гранта Института Прикладной Метафизики
№ 2026-TODO-$\infty$ <<TODO — название гранта>>.

\section*{Доступность данных}
Данные доступны по обоснованному запросу -- а также по необоснованному, с той же
вероятностью успеха.

\begin{thebibliography}{9}
% TODO: 2–3 вымышленных источника; фамилии авторов образованы от терминов этой статьи.
\bibitem{todo1} TODO, Т.~Т. \textit{TODO.} Труды Института Прикладной Метафизики, б.г.
\bibitem{todo2} TODO, Т.~Т. \textit{TODO.} Неопубликованная рукопись.
\end{thebibliography}

\noindent\textbf{Конфликт интересов.} Не заявлен, хотя, согласно Теореме~N,
его отсутствие строго доказать невозможно.

\end{document}
